Ao considerar os trabalhos de \citeonline{barkat2024smart} e
\citeonline{jawas2018water}, que propõem em seus respectivos estudos a
avaliação de modelos de \gls{ocr} baseados em redes neurais aplicados à leitura
de medidores, este trabalho adotou como indicadores para validar a hipótese
\textit{H1} (comparação entre os modelos em ponto flutuante FP32 e
quantizados em INT8) as métricas:

\paragraph{F1 Score}
Média harmônica entre Precisão e Revocação (\textit{Recall}), calculada a
partir dos Verdadeiros Positivos ($VP$), Falsos Positivos ($FP$) e Falsos
Negativos ($FN$) extraídos da Matriz de Confusão, conforme a
Equação~\eqref{eq:f1_score}:

\begin{equation}
    \label{eq:f1_score}
    \text{F1-Score} = \frac{2 \times VP}{2 \times VP + FP + FN}
\end{equation}

Essa métrica avalia o equilíbrio de classificação entre as dez classes
numéricas (dígitos de 0 a 9), evitando distorções causadas pelo
desbalanceamento natural das leituras (como a predominância do dígito 0 nas
casas mais significativas do hidrômetro).

\paragraph{Acurácia de Leitura Total}
Conforme discutido por \citeonline{barkat2024smart}, a eficácia de modelos de
\textit{Deep Learning} aplicados ao monitoramento deve ser mensurada pela sua
capacidade de generalização frente à verdade de campo. A acurácia global obtida
a partir da Matriz de Confusão ($Acc_{ocr}$) relaciona os acertos e erros
totais segundo a Equação~\eqref{eq:acc_ocr}:

\begin{equation}
    \label{eq:acc_ocr}
    Acc_{ocr} = \frac{VP + VN}{VP + VN + FP + FN}
\end{equation}

Em que $VP$, $VN$, $FP$ e $FN$ representam, respectivamente, os Verdadeiros
Positivos, Verdadeiros Negativos, Falsos Positivos e Falsos Negativos. No
contexto operacional do hidrômetro, apura-se a \textbf{Acurácia de Leitura
Total} como a razão entre as imagens em que toda a sequência numérica do
mostrador foi reconhecida sem nenhum erro de caractere e o total de leituras
avaliadas.

\paragraph{Acurácia de Leitura por Dígito}
Corresponde à proporção de recortes individuais de roletes classificados
corretamente em relação ao total de dígitos avaliados, permitindo identificar
se falhas pontuais estão associadas a dígitos específicos, reflexos na cúpula
de vidro ou distorções de perspectiva \cite{barkat2024smart}.

\subsubsection{Métricas aplicadas à desempenho de hardware}
\label{subsubsec:metricas_hardware}

No que tange ao desempenho computacional do \gls{esp32cam}, baseando-se nos
trabalhos de \citeonline{Sari2022} e \citeonline{Ahmad2021}, torna-se
indispensável o levantamento dos compromissos (\textit{trade-offs}) de recursos
embarcados entre os métodos de processamento por meio das seguintes métricas
(hipótese \textit{H2}):

\paragraph{Média de Utilização de CPU durante o processamento}
Mensura a carga computacional e o tempo de execução ativa demandados dos
núcleos do processador do \gls{esp32cam} durante o ciclo de trabalho,
contrastando o esforço exigido pela inferência local das redes neurais no
interpretador TensorFlow Lite Micro (fluxo \textit{Edge}) frente ao custo de
compressão JPEG e empacotamento serial (fluxo \textit{Cloud}), além da
comparação com o tempo de inferência no servidor via LiteRT.

\paragraph{Comparação de frequência de operação entre os métodos durante processamento}
Avalia a permanência do microcontrolador em sua frequência nominal de operação
ativa de 240~MHz durante o processamento dos dados em cada paradigma,
confrontando o tempo em que a CPU permanece em plena atividade antes de
retornar ao estado de repouso em \textit{deep sleep}.

\paragraph{Pico de utilização de PSRAM durante processamento}
Quantifica o volume máximo de memória externa PSRAM (em bytes/megabytes, frente
ao limite físico de 4~MB da placa AI-Thinker \gls{esp32cam}) alocado
estaticamente pelas \textit{arenas} de tensores intermediários dos modelos
embarcados e pelos \textit{buffers} de imagem ($160 \times 120$ pixels)
durante a execução.